\documentclass[a4paper,11pt]{ltjsarticle}
\usepackage[haranoaji]{luatexja-preset}
\usepackage{amsmath,amssymb}
\usepackage[top=12mm,bottom=10mm,left=14mm,right=14mm]{geometry}
\usepackage{array}
\usepackage{tikz}
\usetikzlibrary{calc}
\pagestyle{empty}
\setlength{\parindent}{0pt}
\newlength{\qcol}\setlength{\qcol}{0.47\textwidth}
\newlength{\qgap}
\newlength{\anscolw}\setlength{\anscolw}{0.30\textwidth}
\newlength{\ansh}
\newcommand{\Q}[2]{\parbox[t]{\qcol}{\raggedright (#1)\hspace{0.6em}$\displaystyle #2$}}
\newcommand{\QR}[4]{\Q{#1}{#2}\hfill\Q{#3}{#4}\par\vspace{\qgap}}
\newcommand{\anscell}[1]{\rule[-\ansdrop]{0pt}{\ansh}(#1)}
\newlength{\ansdrop}

\begin{document}
{\large\bfseries 計算問題　2}\hfill 名前(\underline{\hspace{32mm}})\par\vspace{3mm}
■（1）〜（16）の計算をしなさい。（17）〜（20）は方程式を解きなさい。\par\vspace{3mm}
\setlength{\qgap}{5.4mm}
\QR{1}{15a\div(-3)}{2}{(-3)\times\dfrac{2}{3}b}
\QR{3}{-5+(-10)}{4}{(-48)\div(+6)}
\QR{5}{(5x+3)+(2x-7)}{6}{(a+5)-(-2a-5)}
\QR{7}{(+1)+(-10)-(-5)}{8}{5\div3\times(-9)}
\QR{9}{(-8)\times(-3)}{10}{3+2^2+(-4)}
\QR{11}{1-(-6)\div\left(-\dfrac{3}{4}\right)}{12}{\dfrac{1}{3}(9a+6)-\dfrac{1}{6}(12a-6)}
\QR{13}{\dfrac{x-3}{2}\times6}{14}{(24a-9)\div3}
\QR{15}{(-5)-(-1)}{16}{\dfrac{1}{3}(9x-3)+\dfrac{1}{2}(4x+8)}
\QR{17}{0.2x-0.3=-1}{18}{\dfrac{1}{3}x=\dfrac{1}{4}x+5}
\QR{19}{2(x+3)=3x-5}{20}{-4x-7=-3}
\vspace{1mm}
\Q{21}{\{2a-(a-3b)\}+\{2b-(5a-2b)\}}\par\vspace{3mm}
{\small ２問ミスまではまだセーフ　3問以上ミスはもう一回！}\par
{\small 目標は3分！！！　　　Do Your best}\par
\vspace{2mm}
\setlength{\ansh}{9.5mm}\setlength{\ansdrop}{6mm}%
\begin{center}\begin{tabular}{|p{\anscolw}|p{\anscolw}|p{\anscolw}|}\hline
\anscell{1} & \anscell{2} & \anscell{3} \\\hline
\anscell{4} & \anscell{5} & \anscell{6} \\\hline
\anscell{7} & \anscell{8} & \anscell{9} \\\hline
\anscell{10} & \anscell{11} & \anscell{12} \\\hline
\anscell{13} & \anscell{14} & \anscell{15} \\\hline
\anscell{16} & \anscell{17} & \anscell{18} \\\hline
\anscell{19} & \anscell{20} & \anscell{21} \\\hline
\end{tabular}\end{center}

\end{document}
